\section{Security and Formal Analysis}
\label{sec:formal}

This section states the properties that follow from the mechanisms of
Section~\ref{sec:protocol} under the assumptions of Section~\ref{sec:model},
and reports what bounded model checking establishes about them.

\subsection{Analytical Properties}
\label{subsec:analytical}

\textbf{Security of mode switching.} A transition happens only when a
well-formed health certificate implies it (Lemma~\ref{lem:honest-range},
Section~\ref{subsec:health-certificate}). Byzantine validators below $W/3$
cannot move a numeric aggregate outside the range of the honest observations
in a well-formed certificate, and cannot overturn a Boolean value on which all
represented honest validators agree. When honest observations straddle an FSM
threshold, the proposer's choice of quorum subset can still influence which
side of the threshold the aggregate falls on; hysteresis limits this
influence but does not remove it. The certificate therefore guarantees
honest-range validity. It is not an unbiased estimate of the state of the
peripheral layer, and it carries no information about that state beyond what
the honest validators observed. The second attack surface is griefing by an
adversary that degrades the peripheral view of honest validators
(Section~\ref{sec:discussion}).

\refstepcounter{theorem}\label{thm:withdrawal}\textbf{Theorem~\thetheorem{} (Withdrawal Admission Safety).} \emph{Assume $\operatorname{wt}(B) < W/3$, no Bitcoin
reorganization deeper than $K_{\text{DEEP\_FINALITY}}$, and an honest DA
layer. Assume further that (M1) honest validators reject a block that
contains a withdrawal while the state is \textsc{Sovereign} or
\textsc{Recovering}; (M2) each Bitcoin checkpoint commits the chain through
its height, and with it the replicated state; and (M3)
\textsc{Recovering}$\to$\textsc{Anchored} requires a recovery proof that the
chain has verified, covering the \textsc{Sovereign} interval except at most
$T$ trailing blocks (the unproven tail, \texttt{MaxUnprovenTailBlocks}; the heuristic that can stand in for the proof is disabled), and a DA backfill certificate for the whole
interval. If a withdrawal commits in block $B_h$, then
$P_{\text{agree}}$ holds, the state at $h$ is \textsc{Anchored} or
\textsc{Suspicious}, and there are Engram heights $c_s$ and $c_a \geq h - \tau_{\text{DA}}$ such
that $P_{\text{settle}}(c_s)$ and $P_{\text{avail}}(c_a)$ hold, where the
checkpoint that settles $c_s$ lags the certified Bitcoin tip by at most
$\tau_{\text{BTC}}$ Bitcoin blocks. Every height up to $\min(c_s, c_a)$ is
settled and available, including every \textsc{Sovereign} interval that
precedes $B_h$.}

\emph{Proof.} $P_{\text{agree}}$ follows from $\operatorname{wt}(B) < W/3$. By (M1), the state
at $h$ is \textsc{Anchored} or \textsc{Suspicious}. In the transition
function $\delta$, \textsc{Suspicious} is entered only from
\textsc{Anchored}, and \textsc{Sovereign} is left only through
\textsc{Recovering}; so every \textsc{Sovereign} interval before $h$ ended
with a \textsc{Recovering}$\to$\textsc{Anchored} transition, which by (M3)
required a verified recovery proof and a DA backfill certificate. Hence the
interval is hash-linked and available, and all but at most $T$ of its trailing
blocks are covered by the proof; those last blocks are certified by the validator
quorum alone (Section~\ref{subsec:limitations}). In the \textsc{Anchored} and
\textsc{Suspicious} blocks since then, every block carries receipts within
$\tau_{\text{BTC}}$ and $\tau_{\text{DA}}$. The Bitcoin receipt in $B_h$ attests a checkpoint at depth at least
$K_{\text{DEEP\_FINALITY}}$ for some $c_s$; by (M2) it
commits the chain through $c_s$, so rewriting the prefix up to $c_s$
requires a deeper reorganization, which is excluded. The DA receipt in $B_h$
attests data availability up to some $c_a \geq h - \tau_{\text{DA}}$, and the DA layer
is honest. $\square$

\textbf{Scope of Theorem~\ref{thm:withdrawal}.} The theorem is a safety
property of the committed history: it constrains which blocks may contain a
withdrawal. It does not bound what a withdrawal depends on. Settlement is
guaranteed only through $c_s$ and availability only through $c_a$, whereas a
withdrawal may depend on state created inside the exposure window after them,
and the prototype records no dependency height for a withdrawal. Asset-level
safety would add the condition that every withdrawal depends only on state at
heights up to $\min(c_s, c_a)$; the theorem does not establish it. Liveness
comes from relaxing the receipt checks in \textsc{Sovereign}; the property a
bridge would need comes from tying withdrawals to the settled and available
prefix.

The prototype realizes the three assumptions as follows:
\begin{itemize}
  \item (M1) is enforced in the validity predicate and the ante handler.
  \item (M2) holds because each checkpoint payload carries the CometBFT block
  hash at its height $c$. That hash commits the header of $B_c$, hence the
  chain of blocks through $c$ by the hash links, and the header carries the
  application hash of the state before $B_c$. The checkpoint therefore commits
  that state and the transactions of $B_c$, and, since execution is
  deterministic, the state after $B_c$.
  \item (M3) is enforced on the exit edge by the backfill condition, which the
  FSM safety run checks as an action property, and by the proof condition
  (qualified in Section~\ref{subsec:limitations}).
\end{itemize}
The theorem is a pen-and-paper argument over these assumptions, not a
mechanized statement about the implementation.

\subsection{Formal Model}

Model-checking consensus, FSM, and sensors together is intractable at useful
scale. We therefore specify the protocol in TLA+~\cite{lamport2002specifying}
as four layers -- an abstract consensus object, a Tendermint-style consensus
layer, the FSM layer, and a server layer composing them -- and check each
layer with TLC. Apalache~\cite{konnov2019apalache} is used for the certificate
lemma, whose observation vector is symbolic; every row of the ablation study is
checked with TLC.
The target property is: under partial synchrony and an $f$-bounded
adversary, no two honest validators commit conflicting blocks or conflicting
FSM states, and any such conflict implies double-signing evidence.

\textbf{Safety versus liveness.} We follow the separation of concerns of
LiDO~\cite{qiu2024lido}, which verifies Jolteon in Coq; our tools and our
liveness encoding differ. In LiDO, safety is a property of an abstract object
that transfers to a protocol through a refinement mapping. The premise that makes
the abstract object safe is \emph{QuorumOverlap}: any two valid quorums share at
least one honest validator. We call it the linking invariant; with honest
validators' refusal to sign conflicting votes it prevents two conflicting
decisions, and with equal voting power it follows from $n \ge 3f{+}1$. It is a
\emph{safety} premise and does not establish liveness. That the composed server
layer refines the abstract object, i.e.\ that the mapping preserves the invariant
at every step, is the open obligation C1 of Section~\ref{subsec:verification-scope}.

Liveness is argued separately. LiDO obtains it by adding a pacemaker to the
abstract object and stating the liveness requirements as properties of
bounded time segments, under a single GST. Our specification expresses the
same kinds of requirement as TLA+ fairness hypotheses that TLC takes as given.
It models GST as in partial synchrony: a predicate \texttt{AfterGST} on the
model's real time, which never decreases, so that once GST has passed it stays
passed. The predicate concerns the network only; it does not mention the FSM
state or the peripheral layers, so the specified liveness property covers every
FSM state, \textsc{Sovereign} included. Four hypotheses carry the argument: the
consensus layer relies on the first three, which correspond to LiDO's pacemaker
and leader requirements, and the fourth, FSM progress, concerns the FSM layer and
has no counterpart in LiDO. They are:
\begin{itemize}
  \item \emph{Partial synchrony:} behaviours reach GST
  (\texttt{(<>AfterGST) \string~> decided}) and messages between honest
  validators are delivered (weak fairness on \texttt{Push} and \texttt{Pull}).
  \item \emph{Timeout progress:} weak fairness on round-timeout expiry and
  on the $f{+}1$ round-skip (fast-forward) action.
  \item \emph{Proposer progress:} weak fairness on each honest node's
  invocation step, so every honest proposer whose round is reached after
  GST eventually proposes.
  \item \emph{FSM progress:} weak fairness on the FSM transition and on
  environment updates, so a persistently healthy (or unhealthy)
  certificate eventually drives the corresponding transition.
\end{itemize}
Only the FSM-layer liveness properties (circuit-breaker liveness, recovery
initiation, complete re-anchoring) are checked; that layer has no network, so
they need FSM progress alone and no GST. Consensus liveness after GST is specified
on the server layer, but its check did not complete (Appendix~\ref{app:formal}),
so the GST model states our assumptions and is not a verified result for the
composed system.


\subsection{Model-Checked Properties and Attack Resilience}
The model-checked properties and attack resilience are described in Appendix~\ref{app:extended}.

\subsection{Bounded Results and Mechanism Ablation}


\begin{table}[tbp]
\caption{Per-layer, non-refinement verification results (TLC)}
\label{tab:formal-verification}
\centering
\footnotesize
\begin{tabular}{@{}lrrrrrr@{}}
\toprule
\textbf{Layer / property} & \textbf{States gen.} & \textbf{Distinct} & \textbf{Depth} & \textbf{Time} & \textbf{Workers} & \textbf{Viol.} \\
\midrule
Tendermint (safety) & 3{,}668{,}099 & 521{,}958 & 25 & 1\,m\,13\,s & 10 & 0 \\
FSM (safety) & 642{,}091{,}376 & \resFSMStatesExact{} & 16 & 4\,m\,50\,s & 10 & 0 \\
FSM (liveness) & 17{,}977{,}073 & 31{,}105 & 25 & 2\,m\,09\,s & 10 & 0 \\
\bottomrule
\end{tabular}
\\[2pt]
\footnotesize Each row is one completed TLC run of that layer's driver. Depth is
the number of breadth-first levels at which the search terminated with every
state satisfying the state constraint explored, so it is the diameter of the
constrained state graph and not a cap. Bounds: $N{=}4$, $f{=}1$, two rounds and
heights below 10 for Tendermint; the thresholds of Appendix~\ref{app:formal} for
the FSM. Hardware: Apple M1 Pro, 10 cores, 32\,GB RAM.
\end{table}

Table~\ref{tab:formal-verification} reports the per-layer results. Each row is a
run of the layer's own driver at the configuration of
Table~\ref{tab:formal-config}, and none found a violation in its bounded state space.
The FSM safety run, repeated with a different number of workers, gave identical
counts. The runs cover the Tendermint-style and FSM layers. The abstract layer is
exercised only through the open refinement obligation, and the server layer through
that obligation, the certificate lemma and two ablations. The ablation study spans
two layers: the hysteresis, P2P health gate and DA backfill gate belong to the FSM
layer, and the circuit breaker and DA consistency to the server layer, so each
result holds within its own layer's configuration.

\begin{table}[tbp]
\caption{Ablation study: removing one mechanism at a time}
\label{tab:ablation}
\centering
\footnotesize
\setlength{\tabcolsep}{3pt}
\newcommand{\rcc}[1]{\raggedright\arraybackslash #1}
\begin{tabular}{@{}>{\rcc}p{0.14\linewidth}>{\rcc}p{0.20\linewidth}>{\rcc}p{0.19\linewidth}>{\rcc}p{0.06\linewidth}>{\rcc}p{0.08\linewidth}>{\rcc}p{0.21\linewidth}@{}}
\toprule
\textbf{Removed} & \textbf{Targeted threat} & \textbf{Tool, spec} & \textbf{Depth} & \textbf{Time} & \textbf{Violated property} \\
\midrule
Circuit breaker & Withdrawal while locked & TLC, current spec & 5 & 2\,m\,03\,s & \texttt{Sanity\_NeverAttempt} \texttt{WithdrawalWhileLocked} \\
Hysteresis & Premature re-anchoring & TLC, current spec & 9 & 21\,s & \texttt{Sanity\_NoIllegal} \texttt{HysteresisExit} \\
P2P health gate & False recovery under eclipse & TLC, current spec & 5 & 1\,s & \texttt{Sanity\_NoIllegal} \texttt{P2PRecovery} \\
DA consistency & Data withholding & TLC, current spec & 4 & 4\,s & \texttt{Sanity\_Never} \texttt{ProposeWithheldData} \\
DA backfill gate & Exit from \textsc{Recovering} before the interval's data is republished & TLC, current spec & 10 & 49\,s & \texttt{Sanity\_NoIllegal} \texttt{BackfillExit} \\
\bottomrule
\end{tabular}
\end{table}

Table~\ref{tab:ablation} reports the ablation study. Each row removes one mechanism
from the specification of its layer and checks the sanity property that the
mechanism protects. TLC finds a violation in every case, and the counterexample is
the shortest trace to it, since the search is breadth-first. Without the circuit
breaker a withdrawal can be proposed in a locked state; without hysteresis the chain
re-anchors prematurely; without the P2P gate an eclipsed validator forces an illegal
recovery; without DA consistency a withholding proposal is accepted; and without the
backfill gate the chain leaves \textsc{Recovering} before the interval's data is
republished. Two further ablations are not part of this safety table: removing the
$f{+}1$ fast-forward is a liveness ablation and has not completed, and removing the
Byzantine bound of Lemma~\ref{lem:honest-range} is reported with that lemma
(Section~\ref{subsec:formal-invariants}). Single-factor ablation does not test joint
removal, so it does not establish that the set of mechanisms is sufficient.

\subsection{Cross-Layer Refinement and Verification Scope}
\label{subsec:verification-scope}

The evidence is per-layer bounded model checking, not an end-to-end proof. The
Tendermint-style layer (safety) and the FSM layer (safety and liveness) are each
checked standalone, exhaustively within a stated bounded configuration
(Table~\ref{tab:formal-verification}); the ablation study removes one mechanism at
a time (Table~\ref{tab:ablation}); and the certificate lemma is checked
symbolically with Apalache. The abstract and server layers are not checked
standalone. What would transfer the target property to the composed system is the
obligation that the server layer refines the abstract consensus object. It is
specified at three scales (C1: $N{=}4, f{=}1$; C2: deeper rounds; C3:
$N{=}7, f{=}2$) but has not converged at any of them, so the composed system is
\emph{not} formally verified end-to-end; its whole-system evidence in this paper is
empirical (Section~\ref{sec:evaluation}).

Why the refinement of the composed layer is not claimed is explained in Appendix~\ref{app:extended}.

\textbf{What is and is not claimed.} We claim that the Tendermint-style and FSM
layers satisfy the properties of Table~\ref{tab:formal-verification} within the
stated bounds, that each mechanism of Table~\ref{tab:ablation} is needed for the
property it protects when removed on its own, and that the certificate lemma holds
at $N{=}4$, $f{=}1$. We do not claim refinement of the server layer to the
abstract object, safety of the composed system, consensus liveness after GST,
or that the set of mechanisms is sufficient.
